\subsection{A uniform bound for the repeated-quadratic exception}
\label{sec:repeated-quadratic-bound}

All norms in the next argument are taken in the specified quadratic
field, even when the element being normed is rational.

\begin{theorem}\label{thm:non-galois-pu-bound}
Let $\mathcal C$ be a self-dual rank-four pivotal fusion category.
If no Galois conjugate of $\mathcal C$ is pseudo-unitary, then
\[
 \operatorname{FPdim}(\mathcal C)<196.
\]
Consequently its fusion-rule multiplicity is at most $13$.
\end{theorem}
\begin{proof}
Self-duality makes the pivotal structure spherical.
Proposition~\ref{prop:rank4-galois-exception} reduces the question to
$R\otimes\mathbb Q\cong K\times K$ for a real quadratic field $K$.
Write $D=\operatorname{FPdim}(\mathcal C)$ and let $D'$ denote its
conjugate in $K$. Rational $D$ would make $\mathcal C$ weakly integral
and pseudo-unitary, so $D$ is nonrational. The two remaining characters
give formal codegrees $E,E'$, conjugate in $K$, possibly equal if
$E\in\mathbb Q$. All four are positive.

If the spherical dimension character belongs to the Frobenius--Perron
orbit, its category has a pseudo-unitary Galois conjugate. We may
therefore suppose its global dimension belongs to the other orbit,
and conjugate the category so that its global dimension is
$E\geq E'$. Every character value on a basis element has absolute
value at most its Frobenius--Perron dimension; hence every formal
codegree is at most $D$. Thus $E\leq D$, and equality would itself
give pseudo-unitarity. Hence $E<D$.

The formal-codegree theorem supplies a central simple of categorical
dimension $E/D$. In particular, $E/D$ is an algebraic integer in $K$.
Set
\[
 n=\operatorname{Norm}_{K/\mathbb Q}(D)=DD',\qquad
 m=\operatorname{Norm}_{K/\mathbb Q}(E)=EE',\qquad
 k=m/n.
\]
The norm of the nonzero algebraic integer $E/D$ is a positive integer,
so $k\in\mathbb Z_{\geq1}$. The definitions remain valid when
$E=E'$, in which case $m=E^2$.
Put
\[
 A=D+D',\quad B=E+E',\quad p=A/n.
\]
The reciprocal-codegree identity, counting both characters when their
codegrees coincide, gives
\begin{equation}\label{eq:quadratic-reciprocal}
 \frac An+\frac B{kn}=1.
\end{equation}
Consequently
\begin{equation}\label{eq:quadratic-S}
 S:=\frac1{D^2}+\frac1{(D')^2}+\frac1{E^2}+\frac1{(E')^2}
 =p^2+(1-p)^2-\frac{2(1+1/k)}n.
\end{equation}
Ostrik's spherical inequality~\cite[Theorem~3.1.1]{OstrikGlobal} gives
\begin{equation}\label{eq:quadratic-spherical-bound}
 S\leq\frac12+\frac1{2E}.
\end{equation}
For completeness, the squared central dimensions give a positive
measure on twists with total mass $E^2$ and first moment $E$.
The four twist-one constituents of $I(1)$ have squared mass
$h=E^2S$. The remaining mass and first moment are $E^2-h$ and $E-h$,
so $h-E\leq|E-h|\leq E^2-h$. This proves the displayed inequality
without any pseudo-unitarity hypothesis.
Every formal codegree of a nontrivial self-dual fusion ring is strictly
greater than one: it is $1+\sum_{X\ne1}\chi(X)^2$, and the character
cannot vanish on every nonunit basis element because $[X^2:1]=1$.
Thus $D'>1$, so $n>D$; furthermore $E\geq\sqrt{kn}>\sqrt D$.

Suppose for a contradiction that $D\geq196$. First let $k\geq2$.
The function $t+kn/t$ is increasing for $t\geq\sqrt{kn}$, so
$E\leq D$ implies
\[
 B\leq D+kn/D=A+(k-1)n/D.
\]
Using \eqref{eq:quadratic-reciprocal} we obtain
\[
 p\geq\frac{k}{k+1}-\frac{k-1}{(k+1)D}
       \geq\frac23-\frac1{3D}>\frac12.
\]
The last middle expression increases with $k$ because $D>2$.
Equation~\eqref{eq:quadratic-S}, $n>D$, and $k\geq2$ now give
\begin{align*}
 S-\tfrac12
 &>2\left(\frac16-\frac1{3D}\right)^2-\frac3D\\
 &\geq2\left(\frac16-\frac1{588}\right)^2-\frac3{196}
   =\frac1{28}+\frac{589}{172872}>\frac1{28}.
\end{align*}
The expression in the first line is increasing for $D\geq196$.
Since $E>\sqrt D\geq14$, this contradicts
\eqref{eq:quadratic-spherical-bound}.

It remains to treat $k=1$. The algebraic integer $E/D$ is now a
totally positive unit of norm one. Put $r=D/E>1$.
Its two conjugates are $r,r^{-1}$, so $r+r^{-1}$ is an integer
strictly greater than two. Therefore
\[
 r\geq\frac{3+\sqrt5}{2}>\frac52.
\]
Also $E\geq\sqrt n$ implies $r<\sqrt D$. Since $E=D/r$ and
$E'=rD'$, the reciprocal identity gives
\[
 p=\frac r{r+1}-\frac{r-1}{D}.
\]
For fixed $D$ the right-hand side is a concave function of $r$.
Its minimum on $[5/2,\sqrt D]$ occurs at an endpoint. At the first
endpoint it equals $5/7-3/(2D)$, whereas at the second it is
\[
 1-\frac1{\sqrt D+1}-\frac1{\sqrt D}+\frac1D
 >1-\frac2{\sqrt D}\geq\frac67.
\]
Thus $p\geq5/7-3/(2D)>1/2$. Again
\eqref{eq:quadratic-S} gives
\begin{align*}
 S-\tfrac12
 &>2\left(\frac3{14}-\frac3{2D}\right)^2-\frac4D\\
 &\geq2\left(\frac3{14}-\frac3{392}\right)^2-\frac4{196}
    =\frac1{28}+\frac{2249}{76832}>\frac1{28},
\end{align*}
again contradicting \eqref{eq:quadratic-spherical-bound} and $E>14$.
This proves $D<196$ in both cases.

Finally every fusion matrix is real symmetric, and its operator norm
is its Frobenius--Perron eigenvalue $d_X$. Thus
$N_{XY}^{Z}\leq d_X<\sqrt D<14$ for every structure constant.
The multiplicity is an integer and is therefore at most $13$.
\end{proof}

\begin{corollary}\label{cor:pivotal-pu-unrestricted}
A rank-four fusion ring admits a pivotal categorification over
$\mathbb C$ if and only if it admits a pseudo-unitary categorification.
No multiplicity bound is imposed in this statement.
\end{corollary}
\begin{proof}
For non-self-dual rings use the unrestricted non-self-dual theorem.
For a self-dual pivotal category with a pseudo-unitary Galois conjugate
there is nothing to prove. Otherwise
Theorem~\ref{thm:non-galois-pu-bound} puts its fusion ring inside the
verified bounded classification through multiplicity
thirteen. Every positive pivotal ring in that classification has a unitary,
and hence pseudo-unitary, realization. Conversely, a pseudo-unitary
category has its canonical pivotal structure.
\end{proof}

\begin{corollary}\label{cor:simple-pivotal-galois-pu}
Every pivotal rank-four category whose fusion ring has no proper
nontrivial fusion subring has a pseudo-unitary Galois conjugate.
\end{corollary}
\begin{proof}
Such a categorification is self-dual by the unrestricted non-self-dual
classification. A hypothetical exception to Galois pseudo-unitarity
would have multiplicity at most $13$ and rational fusion algebra
$K\times K$. The bounded classification leaves only the rings of
$\operatorname{PSU}(2)_7$ and $\mathcal H_1$ in the absence of a
proper nontrivial fusion subring. Their rational fusion algebras have
degree patterns $1+3$ and $1+1+2$, respectively, rather than $2+2$.
This is a contradiction.
\end{proof}

\begin{remark}
Theorem~\ref{thm:non-galois-pu-bound} bounds the case where the given
category has no pseudo-unitary Galois conjugate. The separate moment
argument of Theorem~\ref{thm:quadratic-product-bound} bounds the
pseudo-unitary case for a product of two quadratic fields as well.
Thus this repeated-field possibility creates
no exception in the unrestricted classification by rational algebra type.
The exception for individual realizations is real: the product of a
Fibonacci category and its Yang--Lee conjugate has global dimension
$E=5$, $D=5\phi^2$, and $n=m=25$, so none of its Galois conjugates
is pseudo-unitary even though its fusion ring has a unitary realization.
\end{remark}
